\section{Horizon charge, extremal area, and arbitrary gauge}\label{sec:horizon-identity}

\subsection{The bifurcation-surface charge}\label{subsec:bifurcation-charge}

Let $\gamma$ be the bifurcation surface of the background Killing field $\xi$. We normalize

\begin{align}
\xi|_\gamma=0,
\qquad
\nabla_\mu\xi_\nu|_\gamma=s_\xi\epsilon_{\mu\nu},
\end{align}

with a fixed binormal orientation. The inner Noether cut is then

\begin{align}
H_\xi^\gamma
=\frac{s_\xi}{2\pi}\frac{A_\gamma}{4G}
\end{align}

on the background. To determine its perturbative extension, keep the surface embedding fixed and write

\begin{align}
g(\lambda)=G+\lambda p+\lambda^2r+O(\lambda^3).
\end{align}

In the project null normal frame $(k,l)$ satisfying $k\cdot l=+1$, the exact local expansion through second order is

\begin{align}
[\lambda]\left(\kappa_{\rm p}^{2}H_\xi^\gamma-s_\xi A_\gamma\right)&=0,\\
[\lambda^2]\left(\kappa_{\rm p}^{2}H_\xi^\gamma-s_\xi A_\gamma\right)
&=\frac{s_\xi}{2}
\int_\gamma\sqrt q\,p_{kk}p_{ll}. \label{eq:fixed-surface-obstruction}
\end{align}

The second-order coefficient $r$ cancels from the difference. Hence, for a smooth two-sided perturbation, a fixed surface, the displayed unit-binormal convention, and an exact background Killing field, the fixed-surface Noether charge equals the fixed-surface area through second order in a Hollands--Wald normal-plane representative satisfying

\begin{align}
h_{kk}|_\gamma=h_{ll}|_\gamma=0. \label{eq:hw-normal}
\end{align}

Equation (\ref{eq:fixed-surface-obstruction}) is the project-convention version of the local obstruction found in the graviton generalized-entropy analysis of Ref. \cite{ColinEllerin:2025gravitons}. It also explains why one cannot identify an arbitrary fixed-coordinate Noether cut with the quadratic area without specifying the normal-plane gauge.

This is a second-order statement. Hollands--Wald gauge is not being used as an all-orders theorem equating a Noether charge with area, and neither smooth extremality nor the gauge condition removes an independent boundary flux in a radiative phase space.

\subsection{Moving extremal surfaces}\label{subsec:moving-extremal-surfaces}

For a displaced surface, metric and embedding variations must be varied together. Let $u$ be arclength along the background geodesic and $V^a$ the first-order normal displacement in a parallel normal frame. In AdS$_3$ the Jacobi operator is

\begin{align}
(JV)_a=(-D_u^2+1)V_a,
\end{align}

while the metric-induced variation of the geodesic curvature is

\begin{align}
\delta_hK_a
=\frac12\nabla_a h_{uu}-D_uh_{ua}.
\end{align}

Linearized extremality is the boundary-value problem

\begin{align}
(JV)_a+\delta_hK_a=0. \label{eq:jacobi}
\end{align}

The metric--displacement and displacement--displacement area terms are

\begin{align}
A^{\rm lin}[h,V]
&=\int_\gamma\mathrm du
\left(\frac12V^a\nabla_a h_{uu}+h_{ua}D_uV^a\right),\\
A^{\rm quad}[G,V]
&=\frac12\int_\gamma\mathrm du
\left(D_uV_aD_uV^a+V_aV^a\right). \label{eq:area-displacement}
\end{align}

For a closed BTZ bifurcation circle, integration by parts produces no endpoint term. For an anchored RT geodesic it produces

\begin{align}
b_{\rm end}
=\left[h_{ua}V^a+\frac12V_aD_uV^a\right]_{\partial\gamma}. \label{eq:area-endpoint}
\end{align}

This term is not discarded; it is included in the joint and anchor regulator analysis of Sections \ref{sec:benchmarks} and \ref{sec:brown-henneaux-theorem}.

Under a metric diffeomorphism $h\mapsto h+\mathcal L_vG$, covariance of the extremality equation gives

\begin{align}
\delta_{\mathcal L_vG}K=Jv_\perp,
\qquad
V\mapsto V-v_\perp .
\end{align}

Thus neither a coordinate location nor a metric perturbation by itself defines the gauge-invariant quadratic area; the pair $(g,\gamma)$ does.

\subsection{Canonical energy and the horizon cocycle}\label{subsec:horizon-cocycle}

The project canonical energy is

\begin{align}
E_{{\rm can,p}}[h]
:=\Omega_{\rm p}[G;h,\mathcal L_\xi h].
\end{align}

Let $h^{\rm ext}=h+\mathcal L_vG$ be a representative satisfying (\ref{eq:hw-normal}) and the linearized extremality condition. The change of canonical energy is algebraically

\[\begin{aligned}
E_{{\rm can,p}}[h+\mathcal L_vG]-E_{{\rm can,p}}[h]
={}&\Omega_{\rm p}
[h+\mathcal L_vG,\mathcal L_{[\xi,v]}G]\\
&-\Omega_{\rm p}[\mathcal L_\xi h,\mathcal L_vG].
\end{aligned}\]

For an arbitrary vector $u$, denote by $\mathbb k^{\rm FA}_{u,\rm p}[\delta g]$ the complete finite-action surface descent of $H_u=X_u\cdot\theta-\alpha_u$. It contains the bulk Iyer--Wald form \cite{Iyer:1994ys}, the timelike-boundary descent, and any required joint or embedding terms. The induced inner-cut cocycle is

\begin{align}
\Upsilon_{\rm p}^{\rm FA}[h,v]
=\mathbb k^{\rm FA}_{[\xi,v],\rm p}
[h+\mathcal L_vG]
-\mathbb k^{\rm FA}_{v,\rm p}[\mathcal L_\xi h]. \label{eq:finite-action-upsilon}
\end{align}

If $v$ is proper at infinity, the constraints vanish, and the regulator flux is controlled, then

\begin{align}
E_{{\rm can,p}}[h+\mathcal L_vG]-E_{{\rm can,p}}[h]
=\int_\gamma\Upsilon_{\rm p}^{\rm FA}[h,v]. \label{eq:cocycle-energy}
\end{align}

The vector $v$ is generally not Killing. Consequently one must use the unsimplified surface descent in (\ref{eq:finite-action-upsilon}); substituting $\xi\mapsto v$ into a formula simplified with the Killing equation is invalid.

Equation (\ref{eq:cocycle-energy}) gives the completed energy

\begin{align}
E_{{\rm can,p}}^{\rm GI}[h]
=E_{{\rm can,p}}[h]
+\int_\gamma\Upsilon_{\rm p}[h,V[h]]
=E_{{\rm can,p}}(P_{\rm HW}h). \label{eq:completed-energy}
\end{align}

It obeys the composition law of a section-change cocycle. In the finite-action convention of this paper,

\begin{align}
\omega_{\rm p}=-\omega_{\rm CELP},
\qquad
\Upsilon_{\rm p}=-\Upsilon_{\rm CELP}.
\end{align}

The sign map changes the separated representatives, not the gauge-invariant combination.

More explicitly, under a further proper section change $h\mapsto h+\mathcal L_wG$ and $v\mapsto v-w$, the integrated cocycle changes by

\begin{align}
\int_\gamma\Upsilon_{\rm p}[h+\mathcal L_wG,v-w]
-\int_\gamma\Upsilon_{\rm p}[h,v]
=E_{{\rm can,p}}[h]-E_{{\rm can,p}}[h+\mathcal L_wG].
\label{eq:upsilon-section-law}
\end{align}

Thus $\Upsilon_{\rm p}$ is section dependent; only (\ref{eq:completed-energy}) is invariant. If $w$ changes the asymptotic sources or frame, or carries a nonzero boundary Hamiltonian, it is not a proper section change and (\ref{eq:upsilon-section-law}) must retain that charge.

\subsection{Conditional finite-cutoff implication}\label{subsec:conditional-implication}

The preceding identities give a useful general statement. Suppose a differentiable regulator family has: (i) a common topology and fixed orientation convention; (ii) an integrable or source-retaining wall polarization; (iii) two-sided smoothness at the limiting bifurcation surface; (iv) controlled joints and embedding variations; (v) a wall-to-corner transgression; (vi) vacuum Einstein dynamics and Brown--Henneaux asymptotics; and (vii) a proper Hollands--Wald representative. If the total regulator flux in (\ref{eq:outer-minus-inner}) vanishes, then

\begin{align}
\delta^2H_\xi^\infty
=\frac{s_\xi}{2\pi}
\delta^2\!\left(\frac{A[g,\gamma]}{4G}\right)
+E_{{\rm can,p}}^{\rm GI}[h]. \label{eq:conditional-horizon-identity}
\end{align}

This statement is conditional for a generic horizon phase space. Sections \ref{sec:benchmarks}--\ref{sec:brown-henneaux-theorem} verify its hypotheses for the linear vacuum Brown--Henneaux sector in AdS$_3$. The higher-dimensional local-Rindler test below shows that smooth extremal gauge alone does not verify hypothesis (ii) or the vanishing-flux premise for a general radiative slab.
